\selectlanguage{american}

\section{Cloud Service}

As described in our technical architecture we want to store the measured data in a database for later analysis.
We decided to use \href{https://www.mongodb.com/de-de}{MongoDB} as the product \href{https://www.mongodb.com/products/platform/atlas-database}{MongoDB Atlas} has a generous free tier with 512 MB storage.
For our cloud platform we decided to use \href{https://azure.microsoft.com/de-de/products/functions/}{Azure Functions} from Microsoft. 
It is a serverless product for executing custom code on custom events (in our case a HTTPS-POST-Request).
When the function is executed the following steps are executed:

\begin{enumerate}
    \item establish a connection to the previously mentioned database.
    \item crate a record indicating that the function was invoked (useful for debugging connections).
    \item after that we decode the received body.
    \item add the current time to the data.
    \item store the body information in the database
    \item return the current time as a response. This is necessary, as the Real-Time Clock drifts depending on the used oscillator crystal.
\end{enumerate}

The code below is written in python and implements the steps discussed above. It is a very, very crude implementation and should only be used in a prototype or evaluation.

\begin{python}[caption={Implementation of a azure function to store incoming data}]
import azure.functions as func
import os
from pymongo import MongoClient
import datetime

MongoDB_URI = os.environ["MONGODB_CONNECTION_STRING"]

app = func.FunctionApp()

@app.route(route="save_to_db", auth_level="anonymous", methods=["POST"])
def save_to_db(req: func.HttpRequest) -> func.HttpResponse:

    client = MongoClient(MongoDB_URI)
    # create db and collection if not exists
    db = client['TheForrest']
    collection = db['Writes']

    # insert current date and time and access ip
    collection.insert_one({"date": datetime.datetime.now(), "ip": req.headers.get('x-forwarded-for')})

    collection = db['TestData']

    try: 
        req_body = req.get_json()
        req_body["server_time"] = datetime.datetime.now()
        collection.insert_one(req_body)
    except ValueError:
        req_text = req.get_body().decode('utf-8')
        collection.insert_one({"data": req_text})
        return func.HttpResponse("Invalid JSON", status_code=400) 
    
    time = datetime.datetime.now(tz=datetime.timezone.utc).replace(microsecond=0).isoformat().replace("+00:00", "Z")

    return func.HttpResponse(time, status_code=200)
\end{python}

The collected data is then visualized in a live \href{https://charts.mongodb.com/charts-project-0-kgbwnwv/public/dashboards/67872274-405b-4a6d-8a79-8702a92f3a80}{Dashboard} (see figure \ref{fig:dashboard_humidity} and \ref{fig:dashboard_battery}). It can be configured to display the the data for a specific time period, by device for humidity and battery level.

\begin{figure}[h!]
    \centering
    \includegraphics[width=1\linewidth]{Bilder/dashbaord.png}
    \caption{Humidity graph of the dashboard}
    \label{fig:dashboard_humidity}
\end{figure}

\begin{figure}[h!]
    \centering
    \includegraphics[width=1\linewidth]{Bilder/dashboard_battery.png}
    \caption{Battery graph of our dashboard}
    \label{fig:dashboard_battery}
\end{figure}

The query used for the dashboard is as follows:

\begin{enumerate}
    \item It checks for the presence of a data array, and the server time.
    \item Sort the remaining entries by the server time (newest first).
    \item Extract the data-points from the data array.
    \item Convert and label the data-points for the dashboard.
\end{enumerate}

\begin{python}[caption={Implementation of a MongoDB Aggregation query for our data}]
from pymongo import MongoClient

client = MongoClient('mongodb+srv://.../')
result = client['TheForrest']['TestData'].aggregate([
    {
        '$match': {
            'server_time': {
                '$exists': True
            }, 
            'data': {
                '$exists': True
            }
        }
    }, {
        '$sort': {
            'server_time': -1
        }
    }, {
        '$unwind': '$data'
    }, {
        '$project': {
            'id': '$data.id', 
            'server_time': 1, 
            'timestamp': {
                '$toDate': {
                    '$multiply': [
                        '$data.t', 1000
                    ]
                }
            }, 
            'humidity': '$data.mV_1', 
            'battery': '$data.mV_2'
        }
    }
])
\end{python}

